\chapter{对称}\label{ch:g4:symmetry}

蝴蝶、脸庞、雪花：大自然偏爱两半
对折后恰好重合的图形。本章借助折纸和方格纸
探索对称。用三角板作图的内容留待
\cref{ch:g6:symmetry}。

\section{折叠与重合}

\begin{definition}[对称轴]\label{def:g4:symmetry:def}
若某直线是图形的\emph{对称轴}\index{对称轴}，
则沿该直线对折图形时，图形的两半
\emph{完全}重合。一个图形可以没有对称轴，
也可以有一条或多条对称轴。
\end{definition}

\begin{omfigure}
\begin{tikzpicture}[scale=0.6]
  % heart-ish symmetric shape: isosceles triangle
  \draw[thick, fill=omDef!12] (0,0) -- (2.4,0) -- (1.2,2.6) -- cycle;
  \draw[omThm, dashed, thick] (1.2,-0.5) -- (1.2,3.0);
  \node[below, font=\small] at (1.2,-0.7) {一条对称轴};
  % rectangle: two axes
  \begin{scope}[xshift=5.4cm]
  \draw[thick, fill=omDef!12] (0,0.3) rectangle (3.2,2.3);
  \draw[omThm, dashed, thick] (1.6,-0.4) -- (1.6,3.0);
  \draw[omThm, dashed, thick] (-0.5,1.3) -- (3.7,1.3);
  \node[below, font=\small] at (1.6,-0.7) {两条对称轴};
  \end{scope}
  % scalene: none
  \begin{scope}[xshift=11.2cm]
  \draw[thick, fill=omDef!12] (0,0) -- (3.0,0.4) -- (0.8,2.2) -- cycle;
  \node[below, font=\small] at (1.4,-0.7) {没有对称轴};
  \end{scope}
\end{tikzpicture}

{\small 沿虚线对折，两半重合。最后一个
三角形无论怎样对折都不能重合，因此没有对称轴。}
\end{omfigure}

\begin{example}[用透明纸检验]\label{ex:g4:symmetry:trace}
描下图形，把透明纸沿假定的对称轴翻过来，
再放回原处：若描下的图形与原图形完全重合，
则该直线确是对称轴。仅凭眼睛往往会判断错误——长方形的
对角线\emph{看起来}像对称轴，但对折检验表明它不是
（不妨自己试一试！）。
\end{example}

\section{在方格纸上补全图形}

\begin{method}[关于方格线的轴对称图形]\label{met:g4:symmetry:grid}
对称轴是方格中的一条竖直（或水平）直线。
\begin{enumerate}
  \item 依次取图形的每个顶点；
  \item 数出它到对称轴的方格数；
  \item 在与对称轴距离\emph{相同}的
        \emph{另一侧}、同一行（或同一列）上画出对应顶点；
  \item 按相同顺序连接对应顶点，并检查：所得
        图形与原图形互为镜像。
\end{enumerate}
\end{method}

\begin{omfigure}
\begin{tikzpicture}[scale=0.55]
  \draw[gray!40, thin] (0,0) grid (12,5);
  \draw[omThm, very thick] (6,-0.4) -- (6,5.4);
  \draw[thick, omDef, fill=omDef!15]
    (1,1) -- (5,1) -- (5,2) -- (3,2) -- (3,4) -- (1,4) -- cycle;
  \draw[thick, omMeth, fill=omMeth!15]
    (11,1) -- (7,1) -- (7,2) -- (9,2) -- (9,4) -- (11,4) -- cycle;
  \draw[dashed] (5,2) -- (7,2);
  \draw[dashed] (3,4) -- (9,4);
\end{tikzpicture}

{\small 沿红色对称轴补全图形：每个顶点都跳到
另一侧距离相同的位置。与对称轴相接的顶点
保持不动。}
\end{omfigure}

\begin{example}[形状相同、大小相同、方向相反]\label{ex:g4:symmetry:properties}
对称所得图形与原图形的边长和角都
完全相同：对称是复制图形，但会把它翻转过来，就像
左手与右手。若原字母是 F，它的镜像
是反向的 F，而不是另一个 F。
\end{example}

\section{常见图形的对称轴}

\begin{example}[数一数对称轴]\label{ex:g4:symmetry:shapes}
通过折叠可知：
\begin{itemize}
  \item 正方形有 $4$ 条对称轴（两条过相对两边的中点，
        两条沿对角线）；
  \item 长方形有 $2$ 条对称轴（只有过相对两边中点的两条！）；
  \item 等腰三角形有 $1$ 条；等边三角形有 $3$ 条；
  \item 圆有无数条对称轴：过圆心的每一条直线
        都是对称轴。
\end{itemize}
\end{example}

\begin{example}[字母中的对称]\label{ex:g4:symmetry:letters}
A、M、T、U、V 有竖直对称轴；B、C、D、E 有水平
对称轴；H、I、O、X 两个方向都有对称轴。由这些字母组成的单词也可以
对称：MUM 有竖直对称轴；OXO 两个方向都对称。
\end{example}

\section{练习}

\begin{exercise}[$\star$]\label{exo:g4:symmetry:1}
描一描、折一折：下列图形中哪些有对称轴——
正方形；非正方形的长方形；不等边三角形（三边
各不相等）；圆？
\end{exercise}

\begin{exercise}[$\star$]\label{exo:g4:symmetry:2}
画出下列图形的对称轴：正方形；长方形；等边
三角形。每个图形各有几条对称轴？
\end{exercise}

\begin{exercise}[$\star$]\label{exo:g4:symmetry:3}
利用\cref{ex:g4:symmetry:trace}中的描图纸检验，
说明一个 $6 \times 4$ 的长方形的对角线\emph{不是}
对称轴。
\end{exercise}

\begin{exercise}[$\star$]\label{exo:g4:symmetry:4}
在方格纸上画一条竖直的对称轴，以及字母 L（高 3 格，
宽 2 格），位于距对称轴 $2$ 格处，作出它的
轴对称图形。
\end{exercise}

\begin{exercise}[$\star$]\label{exo:g4:symmetry:5}
在方格纸上画一条水平的对称轴，并在其上方画一面小旗（一个 $1 \times
1$ 的三角形，挂在 $3$ 格长的旗杆上）。
利用对称，在对称轴下方补全图形。
\end{exercise}

\begin{exercise}[$\star$]\label{exo:g4:symmetry:6}
某图形与其对称轴在点 $P$ 处相接。点 $P$ 的
像在哪里？请结合折纸图加以说明。
\end{exercise}

\begin{exercise}[$\star$]\label{exo:g4:symmetry:7}
将大写字母 A、B、E、F、H、N、O、T 分类：有竖直对称轴吗？
有水平对称轴吗？两个方向都有，还是都没有？
\end{exercise}

\begin{exercise}[$\star$]\label{exo:g4:symmetry:8}
一个三角形的三条边分别为 $5$ cm、$5$ cm 和 $3$ cm。它是否
有对称轴？折叠时哪些边会互换位置？
\end{exercise}

\begin{exercise}[$\star\star$]\label{exo:g4:symmetry:9}
在方格纸上自己设计一个图形，使它\emph{恰有两条}
对称轴，并画出这两条对称轴。（提示：想一想长方形，
或者十字形。）
\end{exercise}

\begin{exercise}[$\star\star$]\label{exo:g4:symmetry:10}
给出一个对称图形的一半：房子的左半部分，
即一面正方形的墙和一半三角形的屋顶，竖直对称轴位于其右边缘。
请描述或画出完整的房子。补全后右半部分的长度
与左半部分的长度相比有什么关系？
\end{exercise}

\begin{exercise}[$\star\star$]\label{exo:g4:symmetry:11}
艾雅认为：“每个有两条边相等的三角形都有对称轴，
且每个有对称轴的三角形都有两条相等的
边。”请用图形检验她的这一双重断言：等腰三角形、
等边三角形、不等边三角形。折叠能支持她的说法吗？
\end{exercise}
